\chapter{The status of the headline results}\label{ch:statusall}

Table~\ref{tab:statusall} gathers, in one place and in the order of the seven parts, the results a referee is likely to check first.
Each row gives the value as it stands in its chapter, the status label it carries there, and the chapter that establishes it; the survey forecasts are in Section~\ref{sec:sp_euclid}. The
labels are those of the key in the front matter: \derived{} follows from stated premises with no parameter adjusted to data;
\calc{} is computed from a derived formula and published constants or measurements; \calibrated{} is set once from reference data and
then held fixed; \measured{} is read from data by the instruments of this book; \observed{} is a published measurement; \prediction{}
names a measurement not yet made; \conjecture{} is a claim that is not yet a calculation; \openprob{} is a known gap; \fitted{} is a parameter adjusted to match the
quantity it is later compared with; \analogy{} is a correspondence that guides a method but is not a derivation; \interp{} is how a set of results
is read, adding no result of its own. A row with two
labels carries both, in the order of its parts. A distance in $\sigma$ divides the difference by the reference measurement's error alone
(Planck $0.54$, SH0ES $1.04$); with both errors in quadrature the three Hubble-constant rows read $-0.28\sigma$, $-0.68\sigma$ and $8.6\sigma$. \calc
The predictions themselves, with their test dates and falsifiers, are listed in Chapter~\ref{ch:predictions}.

{\small\setlength{\tabcolsep}{3pt}\begin{longtable}{@{}>{\raggedright\arraybackslash}p{0.24\textwidth}>{\raggedright\arraybackslash}p{0.25\textwidth}>{\raggedright\arraybackslash}p{0.195\textwidth}>{\raggedright\arraybackslash}p{0.225\textwidth}@{}}
\caption{Status of the headline results across the seven parts.}\label{tab:statusall}\\
\toprule result & value & status & where \\\midrule\endfirsthead
\toprule result & value & status & where \\\midrule\endhead
\bottomrule\endfoot
\multicolumn{4}{@{}l}{\emph{Part~\ref{part:1}: IAM's Law and the virial theorem}}\\
Entropy-area coefficient & $\eta=c^3/(4\hbar G)=1/(4\ell_P^2)$, $\tfrac14=2\pi/8\pi$; one nat per $4\ell_P^2$ & \derived{} (given $G$) & Ch.~\ref{ch:iams_law} \\
The cell's CpG sites against the holographic capacity of a nucleus-sized area & $2.82\times10^{7}$ sites, 52 orders of magnitude below $1.6\times10^{59}$ bits & \calc & App.~\ref{app:saturation} \\
Equation of state of the record term & $w_{\rm info}=-1-1/(3a)$; $-1.33$ today & \derived & Ch.~\ref{ch:iams_law} \\
Drive in thermal units for a black hole, $Mc^2/k_BT_{\rm BH}=2S_{\rm BH}/k_B$ & $2.1\times10^{77}$ at one solar mass & \calc & Ch.~\ref{ch:iams_law} \\
Virial share in hydrogen & $\langle K\rangle/|\langle V\rangle|=13.606/27.211=\tfrac12$ & \derived & Ch.~\ref{ch:virial_law} \\
Smarr share of Schwarzschild horizons, 1 to $6.5\times10^9\,M_\odot$ & $T_HS/Mc^2=0.5000000000$ & \calc & Ch.~\ref{ch:virial_law} \\
Virial ratio of simulated halos & $2K/|W|=1.1$--$1.3$ (1.02--1.17 with surface pressure) & \observed & Ch.~\ref{ch:virial_law} \\
The cosmic horizon at one half & $\beta_m/\Omega_m=1/2$ & \prediction & Ch.~\ref{ch:virial_law} \\
\multicolumn{4}{@{}l}{\emph{Part~\ref{part:2}: cosmology}}\\
Deformed-entropy cosmologies (Barrow, Tsallis) against an added source & their exponent $\Delta$ or $\delta$ is fitted; IAM keeps the area law and fixes $\beta_m$ from $\Omega_m$ & \interp & Ch.~\ref{ch:theory} \\
The cosmological coupling, fixed in every chain & $\beta_m=\Omega_m/2=0.15765$ & \derived & Ch.~\ref{ch:dual} \\
Coupling left free against the local distance ladder and $f\sigma_8$ & $0.155\pm0.029$, the local rate restated; $f\sigma_8$ alone allows $-0.28\le\beta\le0.47$ & \calc & Ch.~\ref{ch:dual} \\
Collapse exponent required by $E(a)=e^{1-1/a}$ & $n=7/2$ & \derived{} (given the form) & Ch.~\ref{ch:theory} \\
Dark-sector equation of state, CPL tangent at $a=1$ & $w_0=-1.062$, $w_a=-0.012$ & \derived & Ch.~\ref{ch:theory} \\
Lensing response & $\Sigma=1$ at every redshift & \calc & Ch.~\ref{ch:theory} \\
Matter-sector coupling today & $\mu_0=-0.136$; $1-\mu=13.62\,\%$ (not the growth deficit) & \calc & Ch.~\ref{ch:surveys} \\
Level~2 fit to Planck 2018, coupling fixed & $\Delta\chi^2=+0.54$ (IAM higher), no added parameter & \measured & Ch.~\ref{ch:level2} \\
Level~1 fits, coupling fixed & $\Delta\chi^2=+0.96$ (Planck), $+0.56$ (+RSD), $+1.73$ (+BAO), $+1.58$ (+Pantheon+) & \measured & Ch.~\ref{ch:latetime} \\
$\sigma_8$, Level~2 (Run~C $\to$ Run~A) & $0.8087\to0.7998$ ($-1.1\,\%$, $-1.51\sigma$) & \measured & Ch.~\ref{ch:level2} \\
$S_8$, Level~2 & $0.830\to0.822$ ($-0.78\sigma$) & \measured & Ch.~\ref{ch:level2} \\
Every other Level~2 parameter & moves by under $0.1\sigma$ & \measured & Ch.~\ref{ch:level2} \\
$\sigma_8$, Level~1, each data combination & $-1.6\,\%$ & \measured & Ch.~\ref{ch:latetime} \\
Amplitude $\mu_0$ left free & medians $+0.030$ to $+0.064$; posterior reaches the $+0.2$ prior edge; prediction inside every 90\,\% interval & \calc & Ch.~\ref{ch:latetime} \\
Photon-sector Hubble constant (Level~2 chain) & $67.16\pm0.47$ km\,s$^{-1}$\,Mpc$^{-1}$; $-0.37\sigma$ from Planck & \measured & Ch.~\ref{ch:dual} \\
Matter-sector Hubble constant, $H_0\sqrt{1+\beta_m}$ & $72.26\pm0.50$; $-0.75\sigma$ from SH0ES & \derived{} (relation), \calc{} (value) & Ch.~\ref{ch:dual} \\
The term placed in the background Friedmann equation & $H_0=61.45\pm0.42$, $10.9\sigma$ from Planck: background form excluded, the large distortion the test was built to show; in IAM the background does not change & \measured & Ch.~\ref{ch:level2} \\
Supernova Hubble-diagram shape (Pantheon+SH0ES, 1590 light curves) & $\beta_{\rm distance}=-0.035$ (68\,\%: $-0.068$ to $0.000$): $\Lambda$CDM geometry & \measured & Ch.~\ref{ch:dsvalidation} \\
Supernova distances with the $\beta$ term applied & $\beta$ on distances excluded, $\Delta\chi^2=+23.6$: distances follow the photon sector, as the dual-sector picture predicts & \calc & Ch.~\ref{ch:dsvalidation} \\
Effective growth index today & $\gamma=0.585$ ($\Lambda$CDM 0.554; measured $0.633^{+0.025}_{-0.024}$) & \calc, \observed & Ch.~\ref{ch:s8trend} \\
$f\sigma_8$ deficit against $\Lambda$CDM, same early amplitude & 4.25\,\% ($z=0$), 2.17\,\% (0.3), 1.35\,\% (0.5), 0.41\,\% (1) & \prediction & Ch.~\ref{ch:surveys} \\
$E_G$ statistic & $+1.8\,\%$ at $z=0.3$, $+3.6\,\%$ today & \calc & Ch.~\ref{ch:surveys} \\
CMB lensing power & $C_L^{\phi\phi}$ lower by 0.08\,\% & \calc & Ch.~\ref{ch:surveys} \\
ISW--galaxy cross-correlation amplitude & $A_{\rm ISW}^{\rm IAM}/A_{\rm ISW}^{\Lambda\rm CDM}\approx1.03$ & \calc & Ch.~\ref{ch:surveys} \\
Survey test of the IAM $\mu(z)$ & the turn-on of the growth ramp across $0.3\lesssim z\lesssim1$; Fisher forecast with Euclid, Planck lensing and DESI & \prediction, \calc & Sec.~\ref{sec:sp_euclid} \\
Siren sample needed to separate the two Hubble constants at $3\sigma$ & $\sigma(H_0)\le1.70$ km\,s$^{-1}$\,Mpc$^{-1}$ & \calc & Ch.~\ref{ch:surveys} \\
Level~2 $\sigma_8$ against lensing surveys & joint KiDS-Legacy + DES Y3 + DESI $0.1\sigma$; KiDS-Legacy $S_8=0.815$, $0.3\sigma$; DES Y3 3$\times$2pt 0.776, $2.3\sigma$ & \observed & Ch.~\ref{ch:sectortension} \\
CPL image of $w_{\rm info}$ against DESI DR2 & $(-\tfrac43,-\tfrac13)$, 7.6--10.2$\sigma$ from DESI in $w_0$; DESI distances are on the photon ruler, so this is not a test & \calc & Ch.~\ref{ch:wzfuture} \\
Asymptotic expansion rates & 55.57 (light), 70.86 (matter) at $H_0=67.16$; no Big Rip & \calc, \derived & Ch.~\ref{ch:wzfuture} \\
Ratio of matter to photon rate & $H_m/H=1.076$ today, 1.275 as $a\to\infty$ & \calc & Ch.~\ref{ch:darkenergy} \\
Writing rate of the record & peak 3.37\,\%\,Gyr$^{-1}$ at $z=1.26$; 2.53\,\%\,Gyr$^{-1}$ now ($H_0=67.16$) & \calc & Ch.~\ref{ch:darkenergy} \\
Baryon fraction relation $\Omega_b/\Omega_m\approx(3/16)\sqrt{\Omega_\Lambda}$ & holds to 0.79\,\%; derivation of 3/16 open & \observed, \openprob & Ch.~\ref{ch:lambda} \\
Baryon-to-photon ratio from the CMB alone & $\eta=(6.113\pm0.037)\times10^{-10}$ & \observed & Ch.~\ref{ch:baryon} \\
Lensing-to-dynamical mass ratio, Level~1 form only & $1/\mu(z)$; whether it applies in virialized clusters is open & \calc, \openprob & Ch.~\ref{ch:lensdyn} \\
\multicolumn{4}{@{}l}{\emph{Part~\ref{part:bh}: black holes and horizons}}\\
Kerr horizons: bits at the horizon temperature & $T_{\rm BH}S_{\rm BH}/Mc^2=\sqrt{1-\chi^2}/2$ & \calc & Ch.~\ref{ch:blackholes} \\
Stefan--Boltzmann luminosity of a horizon against the Hawking luminosity & $P_{\rm SB}/P_{\rm H}=1$ & \calc & Ch.~\ref{ch:blackholes} \\
Black-body evaporation time of one solar mass & $2.1\times10^{67}$\,yr & \calc & Ch.~\ref{ch:blackholes} \\
\multicolumn{4}{@{}l}{\emph{Part~\ref{part:qp}: particles and quantum records}}\\
Bottom-up collapse exponent from halo mass functions & passes $7/2$ at $z\approx3$--4 & \calc & Ch.~\ref{ch:quantumrecords} \\
Koide ratio of the charged leptons & $Q=0.66666446\pm0.00000508$ (PDG 2024; $0.43\sigma$ from $2/3$) & \observed & Ch.~\ref{ch:koide} \\
Generations allowed by mass positivity at $y/x=\sqrt2$ & none for $n\ge4$ at any offset; at the measured offset $\delta=0.2222$ exactly three & \derived & Ch.~\ref{ch:koide} \\
Electron mass from $H_0$, $m_e\propto H_0^{2/5}$ & $-0.02\,\%$ at $H_0=67.36$; one factor identified numerically & \calc & Ch.~\ref{ch:electronmass} \\
Bell value under pointer-basis decoherence & $S_{\max}=2\sqrt{1+c^2}$, never below 2 at finite time & \derived & Ch.~\ref{ch:entanglement} \\
\multicolumn{4}{@{}l}{\emph{Part~\ref{part:3}: qubits and chips}}\\
Active fluctuators that hold the measured quasiparticle fraction & 60, 600, 6,000 for $10^3$--$10^5\,\mu$m$^3$ islands at $x_{\rm qp}=10^{-7}$ & \derived & Ch.~\ref{ch:xqp} \\
Measured quasiparticle background & $10^{-8}$--$10^{-6}$ in the best-isolated devices & \observed & Ch.~\ref{ch:xqp} \\
Quasiparticle-limited $T_1$ of an Al transmon at 5 GHz & $\approx0.24$ ms at $x_{\rm qp}=10^{-7}$ & \calc & Ch.~\ref{ch:reach} \\
Transmon worked example on the qubit gauge & thermal floor at $6.2\times10^{-4}$; the surface-code threshold reads $10$ above the as-built $A=1$ & \calc & Ch.~\ref{ch:ascoreqc} \\
Thermal floor of a qubit record and of one gate & $p_{\rm eq}=1/(1+e^{hf/\kB T})$; $p_{\rm th}=p_{\rm eq}t_g/T_1$ (detailed balance) & \derived & Ch.~\ref{ch:qplatforms} \\
Per-gate thermal floor of a 5\,GHz transmon at its own 35\,mK & $6.2\times10^{-7}$ ($T_1=68\,\mu$s, 40\,ns gate); at $6\times10^{-4}$ on the gauge of a $10^{-3}$ gate & \calc & Chs.~\ref{ch:ascoreqc},~\ref{ch:qplatforms} \\
Slope of the held-record floor with temperature & $d\ln p_{\rm eq}/d\ln T=M(1-p_{\rm eq})$: 6.85 at 35\,mK, 16.0 at 15\,mK (5\,GHz) & \derived, \calc & Ch.~\ref{ch:thermaln} \\
Coherence-against-material optimum & $T_1^*/T_{1,\rm free}=r/(1+r)$, $r=\sqrt{a/b}$; constants unmeasured & \derived{} (model \conjecture) & Ch.~\ref{ch:walls} \\
Drive optimum of a Rydberg gate & $\Omega^*=(a/2b)^{1/3}$ within the two-term model & \derived{} (model \conjecture) & Ch.~\ref{ch:qplatforms} \\
Switching energy per transition on the chip gauge & Ryzen 9 9950X $576$--$593\times k_BT_j\ln2$ (transistor count the least certain input) & \calc & Ch.~\ref{ch:cmos} \\
Landauer floor against the chip's own junction temperature & $3.33\times10^{-21}$\,J at 75\,\textdegree C, 8.6\,\% higher at 105\,\textdegree C & \calc & Ch.~\ref{ch:chipgen} \\
Generation steps a step series can sustain before the floor & $n_{\rm floor}=\ln R/[-\ln(1-s)]$ & \derived & Ch.~\ref{ch:chipgen} \\
DNMT1 discrimination energy against the measured holding energy & Hopfield gap 1.9--4.4 $k_BT$; measured 3.41 $k_BT$ (3.77 with instrument error removed) & \observed, \measured & Ch.~\ref{ch:onegauge} \\
\multicolumn{4}{@{}l}{\emph{Part~\ref{part:4}: the methylome}}\\
Landauer cost per bit at 37~\textdegree C & $2.968\times10^{-21}$ J & \calc & Ch.~\ref{ch:landauer} \\
Drive of one ATP in thermal units & $M=20.94$; 30.2 Landauer bits & \calc & Ch.~\ref{ch:landauer} \\
Holding energy of a methylated site and the copy-error floor & $E_{\rm hold}=3.41\pm0.12\,k_BT$; $\varepsilon_0=0.032$; $\varphi=0.163$ & \measured, \calc & Ch.~\ref{ch:iama} \\
Neutrophil position on IAM-A & $P=1.149$ (3 donors, whole files, 1.140--1.155) & \measured & Ch.~\ref{ch:iama} \\
Neutrophil reference on EPIC (Met-A) & 0.330263 bits; 6 arrays, 6,000 identity sites & \calibrated & Ch.~\ref{ch:meta} \\
Held-out reading of the six reference arrays & standard deviation 0.020 around 1 & \measured & Ch.~\ref{ch:meta} \\
Floor and the full surface & floor $1\times10^{-7}$ (IAM-A and Met-A); full: Met-A 3.03, IAM-A 4.26 & \calc & Ch.~\ref{ch:gauge} \\
Breach and the cancer region on the gauge & to be measured & \openprob & Ch.~\ref{ch:gauge} \\
Copy-error floor at fixed holding energy & $0.78\times$ at 10~\textdegree C; $1.012\times$ at 38.5~\textdegree C & \calc & Ch.~\ref{ch:temperature} \\
Atlas v2 intervals on held-out data & 92.7\,\% coverage & \measured & Ch.~\ref{ch:atlas} \\
Senescent fibroblasts & unmethylated channel cleaned (0.69) & \measured & Ch.~\ref{ch:gauge} \\
\multicolumn{4}{@{}l}{\emph{Part~\ref{part:5}: synthesis}}\\
Mass at which a black hole and the cosmic horizon share a temperature & $M_{\rm eq}=c^3/4GH$: $2.3\times10^{22}\,M_\odot$ today & \calc & Ch.~\ref{ch:theoryinterp} \\
Direction of time as the direction the horizon ledger is written & each entry costs $k_BT_H\ln2$ & \conjecture & Ch.~\ref{ch:time} \\
Irreversibility criterion for a measurement & $Q\ge Q_L=k_BT_D\ln2$ ($0.0179$ eV at room temperature) & \conjecture & Ch.~\ref{ch:measurement} \\
Quantum-eraser fidelity & $F=1-\exp(-Q_L/Q)$; the form is assumed here; its derivation is open & \conjecture & Ch.~\ref{ch:measurement} \\
Gravitational decoherence time, $10^{-12}$ kg silica sphere at 10 mK & $\tau_{\rm IAM}=509$ s, $\propto T^2$; $\tau_{\rm DP}=7.5\,\mu$s & \prediction & Ch.~\ref{ch:gravdec} \\
Activation function from the full $\Lambda$CDM source term & $E(a)\propto\exp(0.66-0.74/a)$, not $\exp(1-1/a)$ & \calc & Ch.~\ref{ch:nonlocal} \\
Records belong in the entropy of the surface on which they are written & the one identification the book adds & \conjecture & Ch.~\ref{ch:synthesis} \\
The CMOS half and the virial half are one quantity & --- & \conjecture & Ch.~\ref{ch:synthesis} \\
Equivalence of inertial and gravitational mass & $\eta\lesssim10^{-15}$ for ordinary materials & \observed & Ch.~\ref{ch:exploratory} \\
\end{longtable}}

Four points of reading follow from the table. First, the fits to data in Part~\ref{part:2} are made with the coupling fixed by the
law, so the chain values are measurements of the fixed model, not fits of a free parameter; the free-amplitude runs are reported as
medians and bounds because their posteriors reach the prior edge. Second, every floor of the processors in Part~\ref{part:3} is the law at the temperature the record sees, so a device reading needs no
fitted constant; what remains open there is how close each loss channel comes to its floor. Third, the cell gauge of
Part~\ref{part:4} is commissioned on its healthy reference; where breach and the cancer region lie on it is still to be measured.
Fourth, ``no free parameter'' has a precise meaning here. No parameter is fitted to the data the law is tested against: $\beta_m$ comes from
$\Omega_m$; the cell's references are read from healthy cells and then held fixed; the chains with $\mu_0$ free add nothing to the model. That is
narrower than ``no free parameter'' without qualification. Counted in full, the results use calibrated quantities (the cell references held fixed after one
measurement), measured ones (the holding energy $E_{\rm hold}$ and the position $P$), and one factor found numerically, $(2\pi)^{3/10}$ in the electron
mass; the $3/16$ of the baryon relation and the $2/\pi$ of the vacuum estimate are taken as given here; their derivation is an open problem. Each is listed with its label in the rows above. \interp

Every row carries its label, so a reader always knows how strongly each result is claimed and where the next test would bite. Pick a row in your field and push on it.
